% TOP_PROBLEM: {"id": 30006956, "title": "Fujita Very Ampleness Conjecture", "category_id": 5, "category": {"id": 5, "name": "algebraic_geometry", "display_name": "Algebraic Geometry", "description": "Geometric objects defined by polynomial equations.", "slug": "algebraic-geometry", "order_index": 5, "created_at": "2026-07-31T15:26:25.671Z"}, "status": "open"}
% FIRST_POSTED: 2026-09-25
% ENTRY_KIND: statement_only
% !TeX program = lualatex
% Definition and sources only; this file is not an LLM solution attempt.
\documentclass[11pt]{article}
\usepackage[margin=25mm]{geometry}
\usepackage{fontspec}
\setmainfont{Latin Modern Roman}
\usepackage{amsmath,amssymb,mathtools}
\usepackage{unicode-math}
\setmathfont{Latin Modern Math}
\usepackage{xurl,hyperref}
\hypersetup{colorlinks=true,urlcolor=blue}
\setcounter{secnumdepth}{0}
\setlength{\parindent}{0pt}
\setlength{\parskip}{5pt}
\setlength{\emergencystretch}{3em}
\begin{document}
\section*{\texorpdfstring{Fujita Very Ampleness Conjecture}{Fujita Very Ampleness Conjecture}}\label{30006956}
\subsection{Definitions and mathematical statement}
Let $X$ be a smooth complex projective variety of dimension $n$, and let $L$ be an ample Cartier divisor. A line bundle is very ample when its complete linear system defines a closed immersion into projective space. With $K_X$ the canonical divisor, Fujita's assertion here is
\[
 \mathcal O_X(K_X+(n+2)L)\quad\text{is very ample}.
\]
In particular this requires enough global sections not only to avoid base points but also to separate distinct points and tangent directions. The integer $n+2$ is fixed uniformly by dimension. The freeness conjecture with $n+1$ is a different, weaker property and is listed separately in Section~\href{https://mehmetmars7.github.io/Erdosproblems-llm-hunter/problem.html?type=open_problems&id=problem.fujita-freeness-conjecture}{363}. Replacing $L$ by a rational divisor without an integrality convention would not define the same line bundle.
\par\smallskip\textit{Formulation and concept references:} \href{https://arxiv.org/abs/2408.06530}{[S1]}.

\subsection{Short English statement}
Does adding n plus two copies of any ample divisor to the canonical divisor always yield a projective embedding of an n-dimensional smooth variety?
\subsection{Sources}
\begin{itemize}
\item[S1]Murayama, Moving Seshadri constants and effective Fujita-type conjectures (2024).
\url{https://arxiv.org/abs/2408.06530}.
\textit{Formulation source.}
\end{itemize}
\end{document}
